% Importado desde: 2026-03-28_Modelo_Minimo_Hibrido_de_Protoespacio.tex
\chapter{Propuesta inicial: --- modelo minimo hibrido de protoespacio discreto --- con subsector Dirac en el infrarrojo}
\label{ch:modelo-minimo-hibrido}

\section{Objetivo}
Despues de comparar \enquote{grafo local} y \enquote{red causal discreta}, el siguiente paso natural es proponer un modelo minimo concreto que combine ambas ideas.

La meta de este capitulo es construir un candidato sobrio:
\begin{displayquote}
\textbf{una estructura discreta con vecindad local, orientacion causal elemental, fibra interna de cuatro componentes y una regla de actualizacion que produzca Dirac \(3+1\) en el infrarrojo.}
\end{displayquote}

No pretendemos todavia demostrar que este sea \emph{el} protoespacio correcto. Solo queremos fijar un modelo minimo de trabajo que:
\begin{enumerate}[label=\arabic*.]
    \item no dependa conceptualmente de Brillouin como punto de partida;
    \item incorpore causalidad en la estructura misma;
    \item y siga siendo compatible con la construccion efectiva de Dirac que ya controlamos.
\end{enumerate}

\section{Los datos del modelo}
La propuesta minima sera
\begin{equation}
\cP_{\mathrm{min}}=(X,\cN,\prec,\cH_{\mathrm{loc}},\cU,\cG_{\mathrm{micro}}).
\label{c28:eq:Pmin}
\end{equation}

\subsection{Conjunto de eventos}
Tomamos
\begin{equation}
X=\mathbb{Z}\times\mathbb{Z}^3.
\end{equation}

Un elemento de \(X\) se escribe como
\begin{equation}
(n,\bm{m}),
\qquad
n\in\mathbb{Z},
\qquad
\bm{m}=(m_x,m_y,m_z)\in\mathbb{Z}^3.
\end{equation}

La interpretacion es:
\begin{enumerate}[label=\arabic*.]
    \item \(n\) etiqueta capas temporales discretas;
    \item \(\bm{m}\) etiqueta vertices de una red cubica en cada capa.
\end{enumerate}

\subsection{Vecindad espacial}
Dentro de cada capa, la vecindad basica es la de primeros vecinos:
\begin{equation}
\cN(\bm{m})=
\{\bm{m}\pm \hat{x},\bm{m}\pm \hat{y},\bm{m}\pm \hat{z}\}.
\end{equation}

\subsection{Relacion causal elemental}
La relacion causal primaria no une puntos arbitrarios, sino eventos en capas consecutivas:
\begin{equation}
(n,\bm{m}) \prec (n+1,\bm{m}')
\end{equation}
si \(\bm{m}'\) esta entre los sitios permitidos por la regla local del paso.

\subsection{Fibra local}
En cada evento vive una fibra de cuatro componentes:
\begin{equation}
\cH_{\mathrm{loc}}\cong \mathbb{C}^2_{\mathrm{quiralidad}}\otimes \mathbb{C}^2_{\mathrm{Weyl}}.
\label{c28:eq:fiber}
\end{equation}

\section{La intuicion geometrica}
El objeto ya no es una red espacial sola, ni una red causal desnuda. Es una red de capas discretas con:
\begin{enumerate}[label=\arabic*.]
    \item conectividad espacial local;
    \item orientacion temporal entre capas;
    \item y una fibra interna que decide como se propaga la amplitud.
\end{enumerate}

\begin{figure}[h!]
\centering
\begin{tikzpicture}[scale=1.0,>=Latex]
  \foreach \x in {0,1.5,3} {
    \foreach \y in {0,1.2,2.4} {
      \fill[blue!60] (\x,\y) circle (1.5pt);
      \fill[blue!60] (\x+5.2,\y+0.8) circle (1.5pt);
    }
  }
  \foreach \x in {0,1.5} {
    \foreach \y in {0,1.2,2.4} {
      \draw[gray!60] (\x,\y) -- (\x+1.5,\y);
      \draw[gray!60] (\x+5.2,\y+0.8) -- (\x+6.7,\y+0.8);
    }
  }
  \foreach \x in {0,1.5,3} {
    \foreach \y in {0,1.2} {
      \draw[gray!60] (\x,\y) -- (\x,\y+1.2);
      \draw[gray!60] (\x+5.2,\y+0.8) -- (\x+5.2,\y+2.0);
    }
  }
  \draw[->,thick,orange!80!black] (1.5,1.2) -- (6.7,2.0);
  \draw[->,thick,orange!80!black] (1.5,1.2) -- (5.2,0.8);
  \draw[->,thick,orange!80!black] (1.5,1.2) -- (6.7,0.8);
  \node at (1.5,-0.45) {\small capa \(n\)};
  \node at (6.7,3.45) {\small capa \(n+1\)};
\end{tikzpicture}
\caption{Imagen minima del modelo hibrido: vecindad local dentro de cada capa y flecha causal entre capas sucesivas.}
\end{figure}

\section{Regla dinamica local}
\subsection{Proyectores internos}
En el bloque Weyl introducimos
\begin{equation}
P_i^\pm=\frac{1}{2}\left(\bbI_2\pm \sigma_i\right),
\qquad i=x,y,z.
\end{equation}

\subsection{Desplazamientos condicionales}
Sea \(T_{\pm \hat{\imath}}\) el operador que desplaza la amplitud una unidad en la direccion \(\pm \hat{\imath}\). Definimos entonces
\begin{equation}
S_i=
T_{+\hat{\imath}}\otimes P_i^+
+
T_{-\hat{\imath}}\otimes P_i^-.
\label{c28:eq:Si}
\end{equation}

Este operador hace algo muy importante:
\begin{enumerate}[label=\arabic*.]
    \item la direccion del salto depende del estado interno;
    \item la propagacion sigue siendo estrictamente local;
    \item la estructura causal del paso queda ya codificada en la propia actualizacion.
\end{enumerate}

\subsection{Bloque Weyl}
El paso Weyl positivo se toma como
\begin{equation}
U_W=S_zS_yS_x.
\label{c28:eq:UW}
\end{equation}

El bloque opuesto se toma como
\begin{equation}
U_{-W}=S_x^{-1}S_y^{-1}S_z^{-1}.
\label{c28:eq:UmW}
\end{equation}

\subsection{Paso Dirac}
En la fibra ampliada de cuatro componentes:
\begin{equation}
U_0=
\begin{pmatrix}
U_W & 0\\
0 & U_{-W}
\end{pmatrix},
\end{equation}
y luego introducimos una mezcla local de masa
\begin{equation}
M=\ee^{-\,\ii \Delta t\,m\,\tau_x\otimes\bbI_2}.
\label{c28:eq:M}
\end{equation}

El paso total es
\begin{equation}
U_D=MU_0.
\label{c28:eq:UD}
\end{equation}

\section{Por que este modelo es realmente hibrido}
No es un grafo puro porque la regla no es solo adyacencia: tiene orientacion temporal elemental. No es una red causal pura porque la dinamica usa de forma esencial la fibra local y los desplazamientos condicionados.

En una frase:
\begin{displayquote}
\textbf{la causalidad viene del orden de capas y del soporte local del paso; la estructura tipo Dirac viene de la fibra interna y de la mezcla entre bloques Weyl.}
\end{displayquote}

\section{Sector infrarrojo}
Si pasamos al espacio de momentos dentro de una capa homogénea, los desplazamientos dan
\begin{equation}
S_i(\vk)=\ee^{-\,\ii a k_i \sigma_i},
\end{equation}
de modo que
\begin{equation}
U_W(\vk)=
\ee^{-\,\ii a k_z \sigma_z}
\ee^{-\,\ii a k_y \sigma_y}
\ee^{-\,\ii a k_x \sigma_x}.
\end{equation}

Para \(a|\vp|\ll 1\),
\begin{equation}
U_W\simeq \bbI-\ii\Delta t\,H_W,
\qquad
H_W=v(\sigma_x p_x+\sigma_y p_y+\sigma_z p_z),
\end{equation}
con
\begin{equation}
v=\frac{a}{\Delta t}.
\end{equation}

Al pasar a cuatro componentes:
\begin{equation}
H_D=
v\,\tau_z\otimes(\sigma_x p_x+\sigma_y p_y+\sigma_z p_z)
+m\,\tau_x\otimes\bbI_2.
\label{c28:eq:HD}
\end{equation}

Luego el mismo modelo hibrido reproduce el subsector Dirac que ya venimos construyendo.

\section{Que gana esta formulacion respecto de Brillouin}
La ventaja conceptual es importante. La estructura primaria ya no es:
\[
\text{toro de Brillouin + cuasienergia}.
\]

Ahora la estructura primaria es:
\[
\text{eventos + vecindad + orden causal + fibra + dinamica local}.
\]

El lenguaje espectral aparece solo \emph{despues}, si queremos analizar:
\begin{enumerate}[label=\arabic*.]
    \item homogeneidad;
    \item dispersion;
    \item subsector infrarrojo;
    \item estabilidad de ciertos modos.
\end{enumerate}

\section{Que simetrias exactas conserva}
Para el modelo minimo ordenado, las simetrias exactas siguen siendo discretas:
\begin{enumerate}[label=\arabic*.]
    \item traslaciones espaciales \(\mathbb{Z}^3\);
    \item traslaciones temporales discretas \(\mathbb{Z}\);
    \item y el subgrupo puntual residual compatible con el orden del paso.
\end{enumerate}

No asumimos desde el comienzo ni isotropia continua exacta ni grupo de Poincare.

\section{Limites actuales del modelo}
Este modelo minimo aun deja abiertas varias cosas:
\begin{enumerate}[label=\arabic*.]
    \item sigue usando una red regular por comodidad;
    \item todavia no demuestra robustez bajo deformaciones no periodicas;
    \item todavia no sabemos si puede generalizarse manteniendo isotropia efectiva sin depender tanto de homogeneidad exacta;
    \item y todavia no hemos probado que capture una geometria comun para varios observables distintos.
\end{enumerate}

\section{Valor del modelo}
Pese a esas limitaciones, el modelo me parece valioso por tres razones:
\begin{enumerate}[label=\arabic*.]
    \item es el primer candidato concreto de protoespacio hibrido del proyecto;
    \item ya contiene causalidad local y estructura espinorial en el mismo objeto;
    \item y permite seguir usando toda la maquinaria Dirac/Weyl que ya desarrollamos, sin volver a una lectura puramente cristalina.
\end{enumerate}

\section{Conclusion}
La propuesta minima que sale de esta etapa es bastante clara:
\begin{displayquote}
\textbf{un buen candidato inicial de protoespacio no es una red espacial sola ni una red causal desnuda, sino una estructura en capas con vecindad local, orientacion causal elemental, fibra interna y una regla de paso local que en el infrarrojo produzca Dirac.}
\end{displayquote}

Eso nos da por fin un objeto de trabajo positivo, no solo una critica a Brillouin.

\section{Siguiente paso natural}
Desde aqui el paso siguiente mas interesante seria uno de estos dos:
\begin{enumerate}[label=\arabic*.]
    \item deformar suavemente este modelo para romper periodicidad exacta y ver que sobrevive del sector Dirac;
    \item o intentar reescribir el mismo modelo en lenguaje de grafo dirigido o red causal local sin referencia inicial a coordenadas cartesianas.
\end{enumerate}

El primero parece mejor para no perder demasiado control tecnico demasiado pronto.
